\chapter{文献と数値の照合}
\label{ch:verify}

\section{文献の実在の照合}

報告書で引いた文献の一覧（\code{照合/文献の一覧.md}）を、\code{LLM/refcheck.py} で照合した。結果は \code{照合/照合の結果.txt} にある。

（照合の集計は、照合が終わってから作る。）


一致が弱かった 20 本は、著者名を加えて Crossref で照合し直し（上の表の右の列）、書籍と業界誌の記事を含めて、すべての実在を確かめた。照合の後に一覧に加えた Hafsi・Vittori は、Crossref の DOI の記録で確かめた（\ref{sec:othermodel} 節）。

\section{Deep Research の答えの誤りと、確かめられなかった記述}

第 II 部を書くときに見つけた、Deep Research の答え（\code{03\_方針と文献調査/ヘッジファンドの手法/deep/}）の誤りと、原典で確かめられなかった記述を並べる。「確かめた」は、論文の要旨（Crossref・arXiv・NBER のページ）または原典の PDF で確かめたことを指す。

{\small
\begin{longtable}{@{}cP{0.3\linewidth}P{0.32\linewidth}P{0.24\linewidth}@{}}
\toprule
答え & Deep Research の記述 & 照合の結果 & 報告書での扱い \\
\midrule
\endfirsthead
\toprule
答え & Deep Research の記述 & 照合の結果 & 報告書での扱い \\
\midrule
\endhead
\bottomrule
\endfoot
\multicolumn{4}{@{}l}{\textbf{誤り}} \\
1 & Khandani・Lo（2007, 2012）によれば、2007 年 8 月に「標準偏差の 25 倍」の変動が数日続いた & 「25 標準偏差」は、当時の Goldman Sachs の幹部の発言として知られるもので、論文の結論ではない & 逸話として、出所を正して書いた（第 \ref{ch:statarb} 章） \\
1 & Fama・French（2015）の 5 ファクターが「リターンの分散の 95\% を説明」 & 要旨（SSRN の版）にこの値はない & 使わない \\
1 & Menkveld（2013）：1 銘柄 1 日 1,397 回の取引、1 回 0.88 ユーロ、シャープレシオ 9.35 & 要旨にない。要旨は「シャープレシオは資本のコストの仮定に強く依存する」と書く & 使わない。要旨の値（取引への参加率 8.1\%・64.4\%、約 5 分の 4 が受け身）を書いた \\
2 & Lopez-Lira・Tang：GPT-4 が「約 90\% の精度」 & 要旨では、取引できない最初の反応についての、ポートフォリオの日の的中率 & 区別を明記して書いた（第 \ref{ch:text} 章） \\
2 & 文献の一覧の番号と URL & 番号が本文の引用と対応せず、URL の多くが解説のブログ・Wikipedia・別の論文を指す & 一覧は使わず、題名で照合し直した \\
3 & Yang ほか（2022）：データの期間の中央が 2015 年以降の研究では平均のアルファが負 & 2022 年 7 月の版の要旨と本文の検索で見つからない（平均の月 0.36\%、74 研究・1,019 推定値は一致） & 一致した値だけを書いた \\
3 & Jacobs・Müller（2020）：231 のアノマリー & 要旨は 241 & 241 と書いた \\
3 & Sun・Wang・Zheng（2016）の題名「Only Winners in Hard Times Sleep Well: \ldots Market Market Environments」 & 正しくは「Only Winners in Tough Times Repeat: Hedge Fund Performance Persistence over Different Market Conditions」 & 正しい題名で引いた \\
3 & Jorion（2000）・Pástor ほか（2020）の URL & それぞれ Amaranth の論文、別の論文（量的・裁量の運用の規模の不経済）を指す & 題名で照合し直した \\
3 & Ben-David ほか（2023）NBER & NBER の WP は 2020 年（w27454） & 2020 年として引いた \\
\multicolumn{4}{@{}l}{\textbf{原典と一致した数値}} \\
1 & Greenwood・Sammon：指数の効果は 1990 年代 7.6\% → 最近 0.8\% & NBER の要旨と一致（1980 年代は 3.4\%） & 書いた（第 \ref{ch:event} 章） \\
1 & Womack（1996）：買い +2.4\%、売り $-9.1\%$ & 要旨と一致 & 書いた \\
1 & Mitchell・Pulvino（2001）：合併裁定の超過リターン年 4\%、プットの売りに似る & 要旨と一致 & 書いた \\
2 & Korajczyk・Sadka（2004）：容量 約 50 億ドル & 要旨（1999 年 12 月の時価総額の水準）と一致 & 書いた（第 \ref{ch:pcresearch} 章） \\
2 & Hafsi・Vittori の DQN による執行 & 題名・著者が Crossref で一致（2026 年の会議録） & 内容は答えの要約として書いた \\
3 & Ben-David ほか：実効の成功報酬 約 50\%、投資家の取り分 36 セント & NBER の要旨と一致 & 書いた（第 \ref{ch:hfperf} 章） \\
3 & McLean・Pontiff（2016）：標本外で 26\%、公表の後に 58\% 低下 & 分野でよく知られた値（要旨） & 書いた \\
3 & Yang ほか（2022）：74 研究、1,019 推定値、平均 月 0.36\% & 原典の PDF と一致 & 書いた \\
\multicolumn{4}{@{}l}{\textbf{確かめられなかった、または学術論文でない出典}} \\
1 & ペアトレードの年率 8\%・最大ドローダウン 12\%、機械学習のシグナルの減衰 3〜6 か月 & 出典は業界の統計のサイト（Gitnux） & 使わない \\
1 & 個人のファクター投資の超過リターン 年 3.27\% & 出典は修士論文 & 使わない \\
1 & 転換社債の裁定 年 9.40\%（1994〜2004 年） & 出典は大学の論文の抜粋で、著者が不明 & 使わない \\
1 & マクロのファンドのシャープレシオ 1.0・0.8・0.9 & 出典を確かめられない & 使わない \\
1 & Carr・Wu（2009）Table 5 の $-0.594$、売りの情報レシオ 3 超 & 原典で確かめていない & 数値は使わず、結論だけを書いた \\
1 & Bernard・Thomas（1989）Table 2：60 日で 4.2\% & 原典で確かめていない & 数値は使わない \\
2 & BloombergGPT の学習費用 約 267 万ドル、FinGPT 約 65 ドル & 原典で確かめていない & 使わない \\
3 & 2008 年のヘッジファンドの損失「総資産の 7\%、自己資本の 140\%」 & 出典を確かめられない & 使わない \\
3 & LTCM の想定の変動率 8.1\% → 13.6\%・19.2\% & 出典を確かめられない & 使わない \\
3 & ボルマゲドン：「ベータ約 4.0」「VIX 先物 約 93,000 枚」 & 出典は報道・解説 & 使わない \\
\end{longtable}
}

\section{照合の限界}

\begin{itemize}
\item \code{refcheck.py} の題名の照合は、文献の実在と、題名・著者・年の大筋を確かめるもので、本文の数値や主張が論文の内容と合っているかは確かめない。数値は、上の表のように個別に要旨や原典で確かめた。
\item 照合の一致の度合いが高くても、別の論文に一致していることがある（例：Harvey・Liu・Zhu（2016）の題名「\ldots and the Cross-Section of Expected Returns」は、題名が短く、別の論文に一致した）。そうしたものは自分で確かめた。
\item 年は、作業論文の版（SSRN・NBER）の年で一致することがある。報告書の文献の一覧には、主に学術誌に載った年を書いた。
\end{itemize}

\section{別のモデルによる点検}
\label{sec:othermodel}

書き上げた本文の全章を Gemini（gemini-3.1-pro-preview）に 1 回だけ送り、事実の誤り（著者・年・結果の取り違え、数値、定義、歴史）を探させた（\code{照合/gemini\_査読.md}）。指摘は 1 件だけで、「Hafsi・Vittori の論文は実在しない生成 AI の捏造だ」というものだった。しかし、Crossref の DOI（10.1109/acdsa67686.2026.11467851）の記録で、IEEE の 2026 年 2 月の会議録の論文（著者 Yadh Hafsi・Edoardo Vittori、9 ページ）として実在を確かめた。Gemini の知識より新しい文献を「架空」と判定した誤りと考えられる。LLM による照合は、新しい文献の実在の判定には使えない。
